% ===================================================================
%  పట్టిక సంఖ్య 5 — ఇల్లు మరియు దాని నిర్మాణం.
%
%  Traduction, faite en 2026, en telougou d'aujourd'hui, sur l'ido de
%  texto/io. Regles de la colonne : voir 10-pattika-01.tex. Le
%  tableau est un DIALOGUE, et les deux volumes composent le nom du
%  parleur en petites capitales : on garde \textsc{}. Les renvois du
%  plan sont des LETTRES, et l'on ecrit « (1) » pour la salle de bain
%  la ou l'imprimeur l'a compose ainsi au lieu de (l).
%
%  PREMIER TABLEAU ECRIT AVEC LA REGLE DU TABLEAU 4 EN MAIN. Avant
%  de rediger, on a releve dans l'ido toutes les PAIRES DE RENVOIS EN
%  RAPPORT MODIFICATEUR-TETE — « les trois marches DU PERRON », « le
%  seuil DE LA PORTE », « la veranda A PORTE VITREE », « la chambre A
%  ALCOVE », « le decorateur SUR L'ECHELLE », « la scie POUR SCIER LE
%  BOIS », « les ferrures POUR L'ELECTRICIEN SUR LE TOIT DE LA
%  REMISE » —, parce que le telougou place tout modificateur avant sa
%  tete et les sortirait donc a l'envers. Dix paires reperees, dix
%  tetes posees d'abord. C'est la premiere fois de la colonne qu'on
%  travaille sur une liste plutot qu'a l'instinct.
%
%  ET LE RESULTAT EST NET : ZERO INVERSION AU PREMIER JET, sur le
%  tableau LE PLUS LONG du livret — 68 blocs, plus de cent vingt
%  renvois. Les quatre tableaux precedents en avaient coute 2, 0, 9
%  et 7 ; celui-ci, plus gros que les quatre reunis, n'en coute
%  aucune. La regle du tableau 4 n'etait donc pas une explication
%  apres coup : elle se verifie a l'avance, et elle suffit.
%
%  ET LE TELOUGOU REPOND A LA NOTE (*) COMME LE MARATHI Y AVAIT
%  REPONDU. La note passe six lignes a expliquer que l'ido n'a pas de
%  racine pour la solive et qu'il faudrait en adopter une :
%  « Balk-o, teknikala radiko til nun ne adoptita ». Le marathi
%  disait వాసా ; le telougou dit వాసాలు. Deux langues indiennes qui
%  ne sont meme pas de la meme famille, un seul mot ordinaire pour ce
%  que sept langues d'Europe cherchaient a nommer.
%
%  LES METIERS PORTENT LEUR NOM DE CASTE, comme en marathe : తాపీ
%  మేస్త్రి le macon — celui de la truelle —, వడ్రంగి le charpentier,
%  కమ్మరి le forgeron, వడ్డెరలు les tailleurs de pierre, తోటమాలి le
%  jardinier, రంగులవాడు le peintre, కూలీ le manœuvre.
%
%  LA MAISON ELLE-MEME EST TELOUGOUE JUSQU'AU SEUIL : అరుగు le
%  perron — la banquette de pierre devant la porte —, గడప le seuil,
%  నడవ le couloir, సావిడి le salon, వంటగది la cuisine, పడక గది la
%  chambre, మిద్దె la terrasse, నేలమాళిగ la cave, పెరడు la cour,
%  మడులు les planches du jardin, గడియ le verrou, పునాది les
%  fondations, గార le mortier, పరంజా l'echafaudage, కొలిమి la forge,
%  దాగలి l'enclume, పట్టకారు les tenailles, ఉలి le ciseau.
%
%  ET LE MEME TROU QU'EN MARATHE : l'ido separe le karpentisto, qui
%  monte la charpente, du menuzisto, qui fait les portes et les
%  parquets ; le telougou dit వడ్రంగి pour les deux. On l'ecrit deux
%  fois plutot que d'inventer un mot.
%
%  ET UN OUTIL QUE JE NE SAIS PAS NOMMER : la lime (98). Le telougou
%  a un mot pour chaque outil de ce tableau sauf celui-la, et je n'ai
%  pas de forme sure. On decrit — అరగదీసే పనిముట్టు — et le trou est
%  note ici plutot que masque, comme l'hirondelle du tableau 3.
% ===================================================================

%%K t05-apar-1 apar
\VUsaut{6.0mm}
\VUtitre{15.0pt}{60}{పట్టిక సంఖ్య 5}
\VUsaut{1.4mm}
\VUtitre{11.4pt}{0}{ఇల్లు మరియు దాని నిర్మాణం.}
\VUsaut{2.2mm}

%%K t05-tit-1 sub
\VUpk{10.2pt}{40}{మంచి కలయిక.}
\VUsaut{3.0mm}

%%K t05-01-1 p
\textsc{జాన్}. --- మామయ్యా, \VUgras{ఇల్లు} ఎలా కడతారో తెలుసుకోవాలని నాకు చాలా ఉంది.

%%K t05-01-2 p
\textsc{మామయ్య} \textsuperscript{(80)}. --- ఇది చాలా సులువు నాయనా; నాతో రా, శ్రీ బెర్నార్డ్
\textsuperscript{(79)} కట్టిస్తున్న, నేను ఉండబోయే ఇల్లు నీకు చూపిస్తాను.

%%K t05-01-3 p
\textsc{జాన్}. --- మరి ఈ ఇంటి \VUgras{నక్షా} \textsuperscript{(76)} ఎవరు గీశారు?

%%K t05-01-4 p
\textsc{మామయ్య}. --- \VUgras{వాస్తుశిల్పి} \textsuperscript{(75)}; ఆయన చాలా స్పష్టమైన చిత్రం
ఇచ్చారు, అది ఆమోదం పొందింది, ఆ తర్వాత \VUgras{గుత్తేదారు} \textsuperscript{(77)} పని
మొదలుపెట్టారు.

%%K t05-01-5 p
\textsc{జాన్}. --- మరి గుత్తేదారు ఎవరు?

%%K t05-01-6 p
\textsc{మామయ్య}. --- శ్రీ బ్లాంకో, నీ స్నేహితుడు జేకబ్ తండ్రి. వాస్తుశిల్పి శ్రీ రూఫోతో కలిసి
ఆయన కూలీలకు పని చెబుతారు.

%%K t05-01-7 p
అదిగో! శ్రీ బ్లాంకో తన \VUgras{కొలబద్దతో} \textsuperscript{(78)}, శ్రీ రూఫో తన నక్షాతో సరిగ్గా
సమయానికి వస్తున్నారు.

%%K t05-01-8 p
నమస్కారం, పెద్దలారా. ఎలా ఉన్నారు?

%%K t05-01-9 p
\textsc{శ్రీ రూఫో}. --- చాలా బాగున్నాం, ధన్యవాదాలు. నమస్కారం జాన్, ఈ పిల్లాడు ఎంత పెరిగాడు!

%%K t05-01-10 p
\textsc{మామయ్య}. --- అవును నిజంగా, ఇప్పుడు వీడు చిన్న మనిషే. చాలా గంభీరంగా ఉంటాడు; కనిపించే
ప్రతిదాని గురించీ వీడికి చెప్పాల్సిందే. ఈ రోజు ఇల్లు ఎలా కడతారో తెలుసుకోవాలని వీడికి ఉంది.

%%K t05-tit-2 sub
\VUpk{10.2pt}{40}{కూలీలు.}
\VUsaut{3.0mm}

%%K t05-01-11 p
\textsc{శ్రీ బ్లాంకో}. --- సరే, నాతో రా నాయనా, నేను వెంటనే నీకు వివరిస్తాను, ఎందుకంటే
నేర్చుకోవాలనుకునే పిల్లలంటే నాకు చాలా ఇష్టం.

%%K t05-01-12 p
చేతిలో \VUgras{పార} \textsuperscript{(82)} పట్టుకున్న ఆ కూలీ కనిపిస్తున్నాడా? అతను
\VUgras{మట్టిపని కూలీ} \textsuperscript{(81)}. నీటి గొట్టాలు వేసేందుకు అతను \VUgras{కందకం}
\textsuperscript{(46)} తవ్వుతున్నాడు. ఇల్లు నిలబడే చోటి నేలను కూడా అతనే తవ్వాడు, తాపీ మేస్త్రి
నాలుగు \VUgras{గోడలు} లేపేందుకు. ఆ గోడల నేలకింది భాగమే \VUgras{పునాది} \textsuperscript{(2)}.
పునాదుల మధ్య ఖాళీ \VUgras{నేలమాళిగ} \textsuperscript{(3)} అవుతుంది. ఆ నేలమాళిగకు
\VUgras{నేలమాళిగ కిటికీ} \textsuperscript{(5)} అనే ఒక చిన్న కిటికీ, ఒక \VUgras{తలుపు}
\textsuperscript{(4)} మరియు ఒక మెట్లదారి ఉంటాయి.

%%K t05-01-13 p
\VUgras{తాపీ మేస్త్రి} \textsuperscript{(108)} \VUgras{తలుపులకు} \textsuperscript{(8)} మరియు
\VUgras{కిటికీలకు} \textsuperscript{(17)} ఖాళీలు వదులుతూ గోడలు కడుతూ వచ్చాడు.

%%K t05-01-14 p
\textsc{జాన్}. --- కానీ ఆ తాపీ మేస్త్రి నాకు కనిపించడం లేదే!

%%K t05-01-15 p
శ్రీ \textsc{బ్లాంకో}. --- ఇంట్లో అతని పని ఇప్పుడు అయిపోయింది. అతను \VUgras{గుర్రపుశాల}
\textsuperscript{(54)} దగ్గర తన \VUgras{పరంజాపై} \textsuperscript{(110)} ఉన్నాడు; అతను
\VUgras{బట్టలు ఉతికే గది} \textsuperscript{(56)} గోడ కడుతున్నాడు.

%%K t05-01-16 p
అతని దగ్గర తాపీ, తొట్టి మరియు \VUgras{గుండు తాడు} \textsuperscript{(109)} ఉన్నాయి. కానీ ఈ
పేర్లు నీకు గుర్తుండవు.

%%K t05-01-17 p
\textsc{జాన్}. --- తప్పకుండా గుర్తుంటాయి, ఎందుకంటే నేను వెంటనే మామయ్యను చిన్న తాపీ, తొట్టి
అడుగుతాను, గుండు తాడు నేనే దారంతో చేసుకుంటాను.

%%K t05-tit-3 sub
\VUsaut{3.0mm}
\VUpk{10.2pt}{40}{సామగ్రి.}
\VUsaut{3.0mm}

%%K t05-01-18 p
\textsc{మామయ్య}. --- ఆహా! చాలా బాగుంది. కానీ చెప్పు జాన్, ఇల్లు కట్టేందుకు ఏమేమి కావాలి?

%%K t05-01-19 p
\textsc{జాన్}. --- \VUgras{చెక్కిన రాళ్ళు} \textsuperscript{(59)} మరియు \VUgras{గార}
\textsuperscript{(65)} కావాలి.

%%K t05-01-20 p
\textsc{మామయ్య}. --- సరిగ్గా చెప్పావు. తన \VUgras{బండిపై} \textsuperscript{(136)} కూర్చున్న
పెద్ద టోపీ మనిషిని చూడు. అతను \VUgras{బండివాడు} \textsuperscript{(135)}. రోజూ అతను ఇక్కడికి
\VUgras{రాతి దిమ్మెలు} \textsuperscript{(60)} తెచ్చి పోస్తాడు. తన బండిని పెరట్లో ఖాళీ
చేస్తాడు, \VUgras{వడ్డెరలు} \textsuperscript{(83)} ఆ దిమ్మెలను చెక్కి తమ \VUgras{రాతి రంపంతో}
\textsuperscript{(85)} కోస్తారు. \VUgras{కూలీ} \textsuperscript{(146)} వాటిని తాపీ మేస్త్రి
దగ్గరికి మోసుకెళ్తాడు, మేస్త్రి వాటిని అమరుస్తాడు.

%%K t05-01-21 p
ఈలోగా \VUgras{గార కలిపేవాడు} \textsuperscript{(87)} గార తయారుచేస్తాడు. అతనికి \VUgras{ఇసుక}
\textsuperscript{(62)}, \VUgras{సున్నం} \textsuperscript{(63)} మరియు \VUgras{నీళ్ళు}
\textsuperscript{(64)} కావాలి. సున్నం అతని దగ్గరే \VUgras{సంచిలో} \textsuperscript{(71)} ఉంది,
\VUgras{తాపీ మేస్త్రి సహాయకుడు} \textsuperscript{(111)} అతనికి \VUgras{తోపుడు బండిపై}
\textsuperscript{(112)} ఇసుక తెచ్చిస్తాడు, \VUgras{శిష్యుడు} \textsuperscript{(113)} పంపు (58)
దగ్గర ఉన్నాడు. అతను \VUgras{బకెట్} \textsuperscript{(114)} నింపుతాడు. గార త్వరలో
సిద్ధమవుతుంది. \VUgras{గార మోసేవాడు} \textsuperscript{(107)} దాన్ని తీసుకుని నిచ్చెన మీదుగా
పరంజాపైకి మోసుకెళ్తాడు, అక్కడ తాపీ మేస్త్రి ఇంటి ఆ వైపు పూర్తి చేస్తున్నాడు.

%%K t05-01-22 p
ఇప్పుడు, \VUgras{కప్పు} \textsuperscript{(31)} వేసే కూలీని ఏమంటారు?

%%K t05-01-23 p
\textsc{జాన్}. --- అతను \VUgras{కప్పు వేసేవాడు} \textsuperscript{(118)}.

%%K t05-tit-4 sub
\VUsaut{3.0mm}
\VUpk{10.2pt}{40}{ఇంటి కప్పు.}
\VUsaut{3.0mm}

%%K t05-01-24 p
శ్రీ \textsc{బ్లాంకో}. --- అవును, కానీ అతనికంటే ముందు \VUgras{వడ్రంగి} \textsuperscript{(115)}
వస్తాడు.

%%K t05-01-25 p
వడ్రంగి \VUgras{కప్పు చట్రం} \textsuperscript{(35)} తయారుచేస్తాడు. అతను \VUgras{దూలాలను}
\textsuperscript{(37)} \VUgras{మేకులతో} \textsuperscript{(117)} కలుపుతాడు. పైన ఆ వెడల్పైన
మఖమలు చేతుల కూలీలు కనిపిస్తున్నారే, వాళ్ళు వడ్రంగులు. ఒకడు చిన్న దూలం పట్టుకుని ఉన్నాడు,
ఇంకొకడు తన \VUgras{గొడ్డలితో} \textsuperscript{(116)} మేకు చెక్కుతున్నాడు. \VUgras{పొగగొట్టం}
\textsuperscript{(41)} దగ్గర ఉన్నవాడే కప్పు వేసేవాడు. అతను (*) వాసాలపై పట్టీలు కొడతాడు,
పట్టీలపై \VUgras{స్లేటు రాళ్ళు} \textsuperscript{(66)}. అప్పుడప్పుడూ అతను పెంకులు కూడా
వేస్తాడు.

%%K t05-noto-1 noto
\VUnotes{143.2mm}{%
(*) \VUgras{Balk-o} = జర్మన్ \textit{Balken}, ఇంగ్లిష్ \textit{joist}, ఫ్రెంచ్
\textit{solive}, ఇటాలియన్ \textit{travicello}, రష్యన్ \textit{balka}, స్పానిష్ మరియు
పోర్చుగీసు \textit{viga}. ఇంకా స్వీకరించని, కానీ ఫ్రెంచ్ \textit{solive} అనే మాట అర్థాన్ని
ఇచ్చేందుకు సూచించదగిన ఒక సాంకేతిక ధాతువు. అది \textit{trabo} (ఫ్రెంచ్ \textit{poutre}) కాదు,
చిన్న దూలమూ కాదు. అది చతురస్రంగా చెక్కిన కర్రదూలం, అది చెక్క నేలనూ పైకప్పునూ మోస్తుంది.%
}

%%K t05-01-26 p
గుర్రపుశాల కప్పు \VUgras{పెంకుల కప్పు} \textsuperscript{(55)}, ఇళ్ళది \VUgras{స్లేటు కప్పు}
\textsuperscript{(32)}, వరండాది \VUgras{జింక్ కప్పు} \textsuperscript{(24)}.

%%K t05-01-27 p
\textsc{జాన్}. --- జింక్ కప్పులు కూడా కప్పు వేసేవాడే వేస్తాడా?

%%K t05-01-28 p
శ్రీ \textsc{బ్లాంకో}. --- లేదు, అది \VUgras{జింక్ పనివాడి} \textsuperscript{(121)} లేదా
\VUgras{నీటిగొట్టాల పనివాడి} పని — ఇంటి కప్పులో \VUgras{కప్పు కిటికీ} \textsuperscript{(38)}
బిగించేది అతనే. \VUgras{వాన కాలువలూ} \textsuperscript{(43)} \VUgras{కప్పు కాలువలూ}
\textsuperscript{(42)} కూడా అతనే వేస్తాడు. జింకునూ తగరపు రేకునూ అతను వెల్డ్ చేస్తాడు.
\VUgras{పిడుగు నివారిణిని} \textsuperscript{(40)} మరియు \VUgras{గాలి బాణాన్ని}
\textsuperscript{(39)} \VUgras{కప్పు కొనపై} \textsuperscript{(36)} అతనే బిగించాడు.

%%K t05-01-29 p
\textsc{జాన్}. --- అయితే ఇల్లు పూర్తయినట్టేనా?

%%K t05-tit-5 sub
\VUsaut{3.0mm}
\VUpk{10.2pt}{40}{పనిముట్లు.}
\VUsaut{3.0mm}

%%K t05-01-30 p
శ్రీ \textsc{బ్లాంకో}. --- పూర్తిగా కాదు. \VUgras{సున్నం పూసేవాడు} \textsuperscript{(99)}
\VUgras{ఇటుకలతో} \textsuperscript{(61)} ఇంటిని వేర్వేరు గదులుగా విభజించే గోడలు కడతాడు.
\VUgras{సున్నంతో} \textsuperscript{(70)} \VUgras{పైకప్పులకూ} \textsuperscript{(26)} గోడలకూ
అతను పూత వేస్తాడు. ఇప్పుడు అతను \VUgras{వరండా} \textsuperscript{(33)} పైకప్పు చేస్తున్నాడు.
తాపీ మేస్త్రిలాగే అతని దగ్గరా \VUgras{తాపీ} \textsuperscript{(100)} మరియు \VUgras{తొట్టి}
\textsuperscript{(101)} ఉన్నాయి.

%%K t05-01-31 p
ఆ తర్వాత వడ్రంగి తలుపులు, కిటికీలు, \VUgras{చెక్క నేల} \textsuperscript{(16)} మరియు
\VUgras{కిటికీ రెక్కలు} \textsuperscript{(19)} చేస్తాడు.

%%K t05-01-32 p
ఇప్పుడు అతను పెరట్లో తన \VUgras{పని బల్లపై} \textsuperscript{(90)} పనిచేస్తున్నాడు. అతని దగ్గర
\VUgras{రంపం} \textsuperscript{(94)} ఉంది — చెక్క (69) కోసేందుకు —, \VUgras{రంధా}
\textsuperscript{(91)}, \VUgras{బిగింపు} \textsuperscript{(92)}, \VUgras{ఉలి}
\textsuperscript{(94 bis)}, \VUgras{వృత్తలేఖిని} \textsuperscript{(95)} మరియు చెక్క
\VUgras{సుత్తి} \textsuperscript{(95 bis)} ఉన్నాయి.

%%K t05-01-33 p
\textsc{జాన్}. --- ఆహా! కానీ తాపీ పనికంటే వడ్రంగి పని ఎక్కువ సరదాగా ఉంది. మామయ్యా, నాకు పని
బల్ల, రంపం, రంధా కొనిపెడతావా? ఎందుకంటే నేను త్వరలో మేదోర్‌కు ఒక చిన్న ఇల్లు చేయబోతున్నాను.

%%K t05-01-34 p
\textsc{మామయ్య}. --- సరే నాయనా.

%%K t05-01-35 p
\textsc{జాన్}. --- నాకెంత సంతోషంగా ఉందో! మరి ప్రవేశ ద్వారం దగ్గర నిలబడి ఉన్నది ఎవరు?

%%K t05-tit-6 sub
\VUsaut{3.0mm}
\VUpk{10.2pt}{40}{రంగు వేయడం.}
\VUsaut{3.0mm}

%%K t05-01-36 p
శ్రీ \textsc{బ్లాంకో}. --- అతను \VUgras{రంగులవాడు} \textsuperscript{(102)}. \VUgras{రంగు}
\textsuperscript{(103)} డబ్బాతో, తన \VUgras{కుంచెతో} \textsuperscript{(104)} అతను తన నిచ్చెనపై
నిలబడి ఉన్నాడు. ప్రవేశ ద్వారం చెక్క ఎక్కువకాలం నిలిచేందుకు అతను దానికి రంగు వేస్తున్నాడు.

%%K t05-01-37 p
గదుల్లో కాగితాలు అతికించేదీ అతనే.

%%K t05-01-38 p
\VUgras{అలంకరణ పనివాడు} \textsuperscript{(107)} \VUgras{నిచ్చెనపై} \textsuperscript{(106)}
ఉన్నాడు, \VUgras{బాల్కనీలో} \textsuperscript{(28)}, \VUgras{తెర} \textsuperscript{(29)}
బిగిస్తున్నాడు.

%%K t05-01-39 p
పెద్ద కిటికీ దగ్గర నిలబడిన కూలీ \VUgras{గాజు పనివాడు} \textsuperscript{(96)}. అతను
\VUgras{అద్దాలను} \textsuperscript{(18)} \VUgras{పుట్టీతో} \textsuperscript{(73)} బిగిస్తాడు.
నేలమాళిగ తలుపు దగ్గర మోకాళ్ళపై కూర్చున్న ఆ ఇంకో కూలీ \VUgras{తాళాల పనివాడు}
\textsuperscript{(97)}. అతను \VUgras{తాళం} \textsuperscript{(6)} బిగిస్తున్నాడు. అతని
\VUgras{అరగదీసే పనిముట్టు} \textsuperscript{(98)} దగ్గరే ఉంది.

%%K t05-01-40 p
\textsc{జాన్}. --- ఇనుముపై తన \VUgras{సుత్తితో} \textsuperscript{(84)} కొడుతున్న ఆ కూలీ కూడా
తాళాల పనివాడేనా?

%%K t05-01-41 p
శ్రీ \textsc{బ్లాంకో}. --- మరింత కచ్చితంగా చెప్పాలంటే అతను కమ్మరి (125). అతను
\VUgras{ఇనుప సామాను} \textsuperscript{(67)} చేస్తాడు, \VUgras{విద్యుత్ పనివాడి}
\textsuperscript{(124)} కోసం — అతను \VUgras{బండ్ల షెడ్డు} \textsuperscript{(51)} కప్పుపై
ఉన్నాడు. అతని దగ్గర \VUgras{కొలిమి} \textsuperscript{(126)}, \VUgras{దాగలి}
\textsuperscript{(127)} మరియు \VUgras{పట్టకారు} \textsuperscript{(128)} ఉన్నాయి.

%%K t05-01-42 p
విద్యుత్ పనివాడు దూరవాణి \VUgras{రాగి} \textsuperscript{(44)} \VUgras{తీగ} బిగిస్తున్నాడు.

%%K t05-tit-7 sub
\VUpk{10.2pt}{40}{తోటపని.}
\VUsaut{3.0mm}

%%K t05-01-43 p
\textsc{మామయ్య}. --- \VUgras{పెరటి} \textsuperscript{(45)} చివర, \VUgras{ఇనుప జాలి}
\textsuperscript{(47)} దగ్గర, \VUgras{పెద్ద గేటు} \textsuperscript{(48)} దగ్గర నీకు
\VUgras{తోటమాలి} \textsuperscript{(129)} కనిపిస్తాడు. అతను \VUgras{చిన్న తోట}
\textsuperscript{(49)} సిద్ధం చేస్తున్నాడు. తన \VUgras{గడ్డపారతో} \textsuperscript{(131)} నేల
తవ్వాడు, విత్తనాలు చల్లాడు, \VUgras{దంతెతో} \textsuperscript{(130)} చదును చేశాడు. ఆ తర్వాత తన
\VUgras{నీళ్ళు చల్లే డబ్బాతో} \textsuperscript{(132)} \VUgras{మడులకు} \textsuperscript{(150)}
నీళ్ళు పోస్తాడు, మొక్కలు పెరుగుతూ ఉంటాయి.

%%K t05-01-44 p
\VUgras{గుర్రపు కాపరి} \textsuperscript{(133)} ఇప్పుడే గుర్రాన్ని శుభ్రం చేసి దానికి మేత
పెట్టాడు; ఇప్పుడు అతను \VUgras{బండిని} \textsuperscript{(52)} కడుగుతున్నాడు — స్పాంజితో,
\VUgras{చేతిపంపుతో} \textsuperscript{(134)}, బకెట్టెడు నీళ్ళతో. ఆ తర్వాత దాన్ని మళ్ళీ బండ్ల
షెడ్డులో పెట్టి మందపాటి \VUgras{గడియతో} \textsuperscript{(53)} తలుపు వేస్తాడు.

%%K t05-01-45 p
ఇప్పుడు నువ్వు సంతోషంగా ఉన్నావని నా అనుమానం; నీకు తగినంత వివరణ దొరికింది.

%%K t05-01-46 p
\textsc{జాన్}. --- అవును మామయ్యా, కానీ ఇంటి లోపలి భాగం కూడా ఇష్టంగా చూడాలని ఉంది.

%%K t05-01-47 p
\textsc{మామయ్య}. --- అది రేపు అవుతుంది. శ్రీ రూఫో నీకు నక్షా వివరిస్తారు, ప్రతి గది పేరూ
చెబుతారు.

%%K t05-tit-8 sub
\VUsaut{3.0mm}
\VUpk{10.2pt}{40}{ఇంటి లోపలి అమరిక.}
\VUsaut{2.4mm}

%%K t05-apar-2 apar
{\centering\textit{(నక్షా చూడండి.)}\par}
\VUsaut{2.4mm}

%%K t05-01-48 p
\textsc{వాస్తుశిల్పి}. --- రా జాన్, ఇంటి వేర్వేరు భాగాలు నీకు వెంటనే చూపిస్తాను.

%%K t05-01-49 p
మనం \VUgras{కింది అంతస్తులోకి} \textsuperscript{(7)} వెళ్దాం. ఇదిగో \VUgras{గంట మీట}
\textsuperscript{(11)}. మనం మూడు \VUgras{మెట్లు} \textsuperscript{(14)} ఎక్కుతాం —
\VUgras{అరుగు} \textsuperscript{(9)} మెట్లు —, \VUgras{గడప} \textsuperscript{(10)} దాటుతాం —
\VUgras{ప్రవేశ ద్వారం} \textsuperscript{(8)} గడప — ఇప్పుడు మనం \VUgras{మెట్లదారి}
\textsuperscript{(13)} మొదట్లో ఉన్నాం.

%%K t05-01-50 p
\textsc{జాన్}. --- ఇది నడవ.

%%K t05-01-51 p
\textsc{వాస్తుశిల్పి}. --- కాదు, ఇది \VUgras{ప్రవేశ గది} \textsuperscript{(12)}. \VUgras{నడవ}
\textsuperscript{(a)} చివర ఉంది. కుడివైపు ఇదిగో \VUgras{సావిడి} \textsuperscript{(b)} మరియు
\VUgras{భోజన గది} \textsuperscript{(c)}; భోజన గదికి ఎదురుగా \VUgras{వంటగది}
\textsuperscript{(d)} మరియు \VUgras{నౌకర్ల గది} \textsuperscript{(e)}, ఆ తర్వాత ముందు
\VUgras{రాత గది} \textsuperscript{(f)}. \VUgras{వరండా} \textsuperscript{(23)} సావిడి పక్కన
ఉంది, దానికి \VUgras{గాజు తలుపు} \textsuperscript{(25)} ఉంది.

%%K t05-01-52 p
ఇప్పుడు మనం మెట్లు ఎక్కుదాం. \VUgras{పట్టుకమ్మి} \textsuperscript{(15)} ఆధారంగా పట్టుకో,
మెట్లు లెక్కపెట్టు. అవి పద్నాలుగు. ఇదిగో \VUgras{మెట్ల చప్టా} \textsuperscript{(m)}; మనం
\VUgras{మొదటి అంతస్తులో} \textsuperscript{(27)} ఉన్నాం.

%%K t05-tit-9 sub
\VUsaut{3.0mm}
\VUpk{10.2pt}{40}{మొదటి అంతస్తు.}
\VUsaut{3.0mm}

%%K t05-01-53 p
మెట్లదారికి ఎదురుగా \VUgras{మరుగుదొడ్డి} \textsuperscript{(20)} మరియు \VUgras{ప్రసాధన గది}
\textsuperscript{(j)} ఉన్నాయి. కుడివైపు అందమైన \VUgras{పడక గది} \textsuperscript{(g)}, దాని
రెండు కిటికీలు ఇంటి \VUgras{ముఖభాగంలో} \textsuperscript{(1)}. ఎడమవైపు \VUgras{స్నానపు గది}
\textsuperscript{(1)} మరియు \VUgras{అతిథుల గది} \textsuperscript{(i)}, దానిలో పెద్ద
\VUgras{గూడు} \textsuperscript{(h)}: పడక గది పక్కన \VUgras{పిల్లల గది} \textsuperscript{(k)}
ఉంది. అది విశాలమైన \VUgras{మిద్దెకు} \textsuperscript{(21)} ఆనుకుని ఉంది. ఈ మిద్దె కింద ఇంకో
మరుగుదొడ్డి (20) ఉంది, గోడపై ఇదిగో భోజనానికి మోగే \VUgras{గంట} \textsuperscript{(22)}.

%%K t05-01-54 p
\textsc{జాన్}. --- మనం \VUgras{రెండో అంతస్తుకు} వెళ్దాం.

%%K t05-01-55 p
\textsc{వాస్తుశిల్పి}. --- రెండో అంతస్తు లేదు. కప్పు కింద \VUgras{పైగదులు}
\textsuperscript{(33)} మరియు ధాన్యపు గది (34) మాత్రమే ఉన్నాయి. మన పర్యటన అయిపోయింది, మనం
కిందికి దిగవచ్చు.

%%K t05-01-56 p
\textsc{జాన్}. --- ధన్యవాదాలు అయ్యా, ఇది చాలా ఆసక్తికరంగా ఉంది. నేను చూసినదంతా వెంటనే మా
నాన్నకు చెబుతాను.

\VUsaut{4.0mm}
\VUfilet{22mm}
